\scanpage{36}
% L1–4 | top
graduées anticommutatives, $f^*$ \unsure{homogène|homogènes}. Supposons qu'il existe \deleted{deux} entier $d=d(A)$ tels que $k_A$ soit composé $A\unsure{-|\to}A^d\xrightarrow{\unsure{\bar k_A^{(d)}|k_A^{(d)}|\bar{\kappa}_A^{(d)}}} k$, et de même un $d'=d(B)$. Alors si $x$ est de degré $i$, la forme $\varphi_A(x)$ est nulle en degré $\ne d-i$,

% L5–8 | upper
donc $\langle f^*(y),\varphi_A(x)\rangle=k(f^*(y)x)$ est nul \deleted{\unsure{si}} en degré $\ne d-i$, \deleted{\unsure{donc si $f$|donc si $x$}} de $=$, \inserted{si $x'$ \deleted{\unsure{$=f_*(x)$}}, est hom de degré $i'$,} la forme $\varphi_B(x')$ est nulle en degré $\ne d'-i'$, donc on doit avoir (si $\unsure{f_*(x')|f_*(x)|\varphi_B(x')}\ne0$)
\[
d-i=d'-i'\quad\text{i.e.}\quad i'=i+(d'-d).
\]

% L8–12 | upper middle
Si alors on pose $f_*(x)=\sum_j x'_j$, on voit que l'effet de la forme $\varphi_{\unsure{B|A}}(f_*(x))$ sur $B^{d'-j}$ est égal à $\varphi_B(x'_j)$, donc est nul si $\deleted{j}\inserted{d'-j}\ne d-i$, donc dans ce cas $x'_j=0$. Donc

$f_*$ \uline{augmente les degrés de} $d'\unsure{-|+}d$.

% L13–18 | middle
Nous travaillons dans la catégorie \inserted{bivariante} des Algèbres graduées anticommutatives, munies d'un degré $d$, le $f_*:A\to B$ étant assujetti à être de degré $d(B)-d(A)$, \inserted{et $k_A$ \deleted{\unsure{de a|se}} factorisant par $A^{d(A)}$.} La sous catégorie pleine des $A$ tels que $d(A)$ pair est commode

% L19–24 | lower
[et il y a lieu parfois de se borner aux Algèbres $A$ qui sont à degrés \deleted{\unsure{entre}} dans $[0,d(A)]$]. Il n'est pas \deleted{\unsure{utile}} \inserted{commode} de se borner aux ``Algèbres de Poincaré'', associé\unsure{es} comme dessus à des formes sur $A$, car les constructions universelles qu'on va faire en feraient de \unsure{tte} façon sortir !
% archive sheet number: 35

